\section{Le cadre PT}
\label{sec:cadre_pt}

\subsection{Le paramètre fondamental}
\label{sec:parametre}

La théorie de la Persistance (PT) prend pour unique axiome arithmétique
le théorème~$\Tone$ \cite{PT_M1}~:
\[
\sper \;=\; \tfrac{1}{2}.
\]
Ce paramètre joue trois rôles distincts dans le corpus PT, qui seront
progressivement identifiés au cours de l'article~:
\begin{itemize}[nosep,leftmargin=*]
\item \textbf{arithmétique}~: $\sper$ est le paramètre de symétrie de la
  transition interdite au premier $p = 3$ du crible
  (théorème~$\Tone$, \cite{PT_M1}) ;
\item \textbf{géométrique}~: $\sper$ est le seuil $\gammap{p} > 1/2$
  d'activité d'un premier dans la cascade
  (théorème~$\Tfive$, \cite{PT_M3}) ;
\item \textbf{spectral}~: $\sper$ est la partie réelle qui apparaît dans
  la transformée de Mellin sur les premiers via le Jacobien de la mesure
  de Haar multiplicative.
\end{itemize}

Une lecture centrale du présent travail est que ces trois manifestations
sont opératoriellement identifiables (\S\ref{sec:localisation_critique}).

\subsection{Point fixe, premiers actifs, bifurcation}
\label{sec:point_fixe}

Le crible PT admet un unique point fixe $\mustar = 15$. Les premiers
arithmétiquement actifs en ce point fixe sont les diviseurs premiers
$\{3, 5, 7\}$ de $\mustar$ ; le premier $p = 2$ joue le rôle distinct de
canal de parité (cf.~\S\ref{sec:dualite_higgs}). À ce point fixe, la
dynamique de persistance se bifurque en deux branches conjuguées $\qplus$
et $\qminus$ caractérisées par
\begin{equation}
\qplus \;=\; \dfrac{13}{15}, \qquad
\qminus \;=\; \exp\!\Big(-\dfrac{1}{15}\Big).
\label{eq:qplusqminus}
\end{equation}
Cette bifurcation est l'objet central du programme zêta. Elle entre dans
toutes les amplitudes propres et toutes les phases cohomologiques
rencontrées plus bas.

\subsection{Mesure de Haar multiplicative et coordonnée de Mellin}
\label{sec:haar}

La cascade de persistance privilégie une seule mesure invariante sur la
demi-droite des premiers~: la mesure de Haar multiplicative
\begin{equation}
du \;=\; \dfrac{dp}{p}, \qquad u \;=\; \log p.
\label{eq:haar_mult}
\end{equation}
Cette mesure est invariante par dilatation $p \mapsto \lambda p$, et c'est
la seule mesure invariante par le flot RG de PT (théorème~$\Ttwo$,
doublement stochastique de la matrice $T_3$, \cite{PT_M1}). La
coordonnée $u = \log p$ sera nommée \emph{coordonnée de Mellin} dans tout
l'article ; elle porte la mesure additivement invariante par translation
$u \mapsto u + c$.

\subsection{Factorisation PT canonique de $\zeta$}
\label{sec:factorisation}

Pour $\Reop(s) > 1$ on définit, à partir des amplitudes de bifurcation,
\begin{equation}
\delta_p^{\pm} \;=\; \dfrac{1 - q_{\pm}^{p}}{p}, \qquad
a_p \;=\; \delta_p^{+}\,(2 - \delta_p^{+}), \qquad
b_p \;=\; \delta_p^{-}\,(2 - \delta_p^{-}),
\label{eq:apbp}
\end{equation}
puis
\begin{equation}
\zetaplus(s) \;=\; \exp\!\Big(\sum_{p} a_p\,p^{-s}\Big), \qquad
\zetaminus(s) \;=\; \exp\!\Big(\sum_{p} b_p\,p^{-s}\Big),
\label{eq:zetapm}
\end{equation}
et le \textbf{facteur résiduel}
\begin{equation}
\Rres(s) \;=\; \dfrac{\zeta(s)}{\zetaplus(s)\,\zetaminus(s)}.
\label{eq:Rres_def}
\end{equation}
On pose enfin $\Ap = a_p + b_p$, qui sera l'amplitude PT effective par
premier dans toute la suite. Asymptotiquement, $\Ap \sim 4/p$ pour
$p \to \infty$.

\paragraph{Écriture canonique en angles d'holonomie.}
Les amplitudes $a_p, b_p$ s'écrivent aussi
\begin{equation}
a_p \;=\; \sinp{p}(\qplus), \qquad b_p \;=\; \sinp{p}(\qminus),
\label{eq:apbp_sinp}
\end{equation}
soit
\begin{equation}
\Ap \;=\; \sinp{p}(\qplus) \,+\, \sinp{p}(\qminus),
\label{eq:Ap_holonomie}
\end{equation}
où $\thetap{p}$ désigne l'angle d'holonomie de la cascade au premier $p$
(monographie ch.~10, structure fine, \cite{PT_Foundations}). L'équivalence
entre les deux écritures vient de $1 - \delta_p^{\pm} = q_{\pm}^{p}/p$,
soit
\[
\delta_p^{\pm}\,(2 - \delta_p^{\pm}) \;=\; 1 - (1 - \delta_p^{\pm})^{2} \;=\; 1 - (q_{\pm}^{p}/p)^{2} \;=\; \sinp{p}(q_{\pm})
\]
en posant $\cos(\thetap{p}, q_{\pm}) = q_{\pm}^{p}/p$.

Le programme zêta cherche un modèle spectral canonique de $\Rres$, sous
l'\textbf{hypothèse conjecturale} de non-annulation de
$\zetaplus\,\zetaminus$ (vérifiée numériquement sur les zones testées) qui
identifie les zéros non triviaux de $\Rres$ et ceux de $\zeta$. Cette
hypothèse sera rappelée à chaque étape critique de l'identification.

